\section{Capítulo~\ref{sec:glutcomputer}. Qué calculan las neuronas \nt{GLUT}}

Las proposiciones lógicas del capítulo son teoremas sobre circuitos de pesos no negativos, deducidos en el propio capítulo; los experimentos confirman el modelo de neurona y que los circuitos construidos según esos teoremas (detector de coincidencias, línea de retardo, grupo autosostenido, cadena) existen realmente en el cerebro.

{\footnotesize
\setlength{\tabcolsep}{3pt}\renewcommand{\arraystretch}{1.15}
\begin{longtable}{>{\raggedright}p{0.5cm}>{\raggedright}p{4.0cm}>{\raggedright}p{5.6cm}>{\raggedright}p{2.3cm}>{\raggedright\arraybackslash}p{1.7cm}}
\toprule
N.º & Afirmación & Experimento y qué se midió & Fuente & Estado \\
\midrule\endfirsthead
\toprule
N.º & Afirmación & Experimento y qué se midió & Fuente & Estado \\
\midrule\endhead
1 & La neurona es un circuito RC con umbral y reinicio~\eqref{eq:glutneuron} & Al soma de neuronas piramidales de corteza de rata (rebanadas) se aplicó una corriente ruidosa parecida al flujo sináptico del cerebro vivo. La dependencia de la frecuencia media respecto de la media y la dispersión de la corriente se describe con la neurona integradora con fugas si se añade una adaptación lenta de la frecuencia. Los intervalos de $\tau$ y de los retardos se dan en el capítulo sin referencia experimental & \cite{Rauch2003} & medido \\
2 & El modelo~\eqref{eq:glutneuron} es una simplificación: las ramas de entrada generan por sí mismas espigas locales & Registro directo en las dendritas de neuronas piramidales de la corteza visual del ratón in vivo: un estímulo visual provoca espigas dendríticas locales dependientes de los receptores NMDA; su supresión ensancha la sintonía de la neurona a la orientación & \cite{Smith2013} & medido \\
3 & La ventana de coincidencia es $\Delta_{\max}=\tau\ln\frac{w}{\theta-w}$~\eqref{eq:window}; con $\theta=1.5w$ vale $0.7\tau$ & Consecuencia del modelo~\eqref{eq:glutneuron}. Hay una medición directa de una ventana de suma estrecha en neuronas con inhibición rápida (capítulo~\ref{sec:gaba}, fila~3 de su tabla) & deducción del capítulo & teoría \\
4 & Una red formada solo por neuronas \nt{GLUT} calcula únicamente funciones monótonas de las entradas (proposición~\ref{prop:monotone}) & Demostración en el capítulo: iteración desde el silencio con un lado derecho no decreciente. Es una proposición sobre el modelo; no hay un experimento que la haya comprobado en una red viva sin inhibición & deducción del capítulo & teoría \\
5 & Los órganos de los sentidos entregan un par de cables: unas células responden al aumento del estímulo y otras a su disminución & Retina: ya en las células bipolares la señal de los conos se divide en canales de encendido y de apagado. Bloqueo farmacológico de las bipolares de encendido en el mono: los canales van separados hasta la corteza; tras el bloqueo empeora la detección del aumento de luz, pero no la de su disminución & \cite{Schiller1992} & medido \\
6 & Con código de dos cables, una red de neuronas \nt{GLUT} calcula cualquier función lógica de forma asíncrona, sin ciclo de reloj común, al precio del doble de neuronas (proposición~\ref{prop:dualrail}, \eqref{eq:dualrail}, fig.~\ref{fig:dualxor}) & El código de dos cables y la detección de la finalización por la propia señal son una técnica estándar de los circuitos asíncronos insensibles a los retardos. El circuito de OR exclusivo de seis neuronas se construye en el capítulo. No hay experimentos que muestren que el cerebro calcula funciones lógicas precisamente con código de dos cables & \cite{Sparso2001} & teoría \\
7 & La mayor diferencia de tiempo de llegada del sonido a los oídos en el ser humano es $\approx0.6\ms$, y el cerebro distingue unos diez microsegundos & Oyentes en un experimento de elección forzada entre dos alternativas: el umbral de discriminación de la diferencia de tiempo (75\% de aciertos) es de $9$\,\textmu s para ruido en la banda de $150$--$1700$\,Hz, $11$\,\textmu s para un tono de $1$\,kHz y $28$\,\textmu s para un clic. El límite de $0.6\ms$ es un cálculo del capítulo, $0.22\,\text{m}/343\,\text{m/s}$ & \cite{KlumppEady1956} & medido \\
8 & En las aves, la diferencia de tiempos se convierte en lugar mediante líneas de retardo y detectores de coincidencias~\eqref{eq:itd} (fig.~\ref{fig:itd}) & Lechuza común: los axones de los dos oídos entran en el núcleo de detectores de coincidencias por lados opuestos; el tiempo de llegada de las espigas a lo largo del axón cambia de forma ordenada con la profundidad, y el rango de retardos a través del núcleo ($160$\,\textmu s en $700$\,\textmu m) coincide con el rango de retardos interaurales & \cite{Carr1990} & medido \\
9 & En los mamíferos no se ha encontrado una fila de retardos; la misma tarea la resuelven detectores de coincidencias con una inhibición de \nt{GLYC} temporizada con precisión & Registro in vivo en la oliva superior medial del jerbo: las respuestas no forman un mapa de direcciones; la aplicación de glicina y de su bloqueador por micropipeta muestra que una inhibición \emph{glicinérgica} temporizada con precisión fija el rango de trabajo de los retardos. Aquí la inhibición proviene de \nt{GLYC}, no de \nt{GABA} & \cite{Brand2002,Grothe2010} & medido \\
10 & Un grupo con realimentación fuerte es un biestable; la actividad autosostenida dura segundos tras el estímulo (fig.~\ref{fig:bistable}) & Corteza prefrontal de mono en una tarea de movimiento ocular diferido hacia un punto memorizado (retardo de $1$--$6$\,s): $87$ de $170$ neuronas relacionadas con la tarea cambian su frecuencia durante el retardo, y en el $79\%$ de ellas esto depende de la dirección del punto memorizado. Que esta actividad la sostenga una realimentación con dos estados estables es teoría & \cite{Funahashi1989,Wang2001} & medido \\
11 & El bit se escribe si el aporte dura más de $\approx0.7\tau$ (fig.~\ref{fig:recurrentgroup}b) & Cálculo de la ecuación $\tau\,\dd r/\dd t=-r+f(wr+I)$ por el método de Euler con $w=1$, $I=0.6$; no hay script en \texttt{models/}. No hay experimento & cálculo del capítulo & modelo \\
12 & La memoria puede guardarse en la facilitación sináptica, casi sin espigas & Teoría: el calcio residual en las terminales mantiene lo escrito cerca de un segundo. Apoyo: en la corteza prefrontal medial del hurón (registro de varias neuronas a la vez) hay una subred de neuronas piramidales conectadas por sinapsis facilitadoras con un efecto residual prolongado. Revisión de observaciones de intervalos “silenciosos” durante el mantenimiento en memoria. Qué forma usa la corteza y cuándo es una pregunta abierta & \cite{Mongillo2008,WangMarkram2006,Stokes2015} & hipótesis \\
13 & La memoria se borra por fatiga: se agota la reserva de neurotransmisor & Registros por pares de neuronas piramidales de la corteza: la respuesta de la sinapsis cae con espigas frecuentes, y la velocidad de la caída la fija la probabilidad de liberación. Que sea precisamente esto lo que apaga un grupo autosostenido no se ha mostrado directamente & \cite{Tsodyks1997} & medido \\
14 & Una cadena de grupos es un temporizador disparado por un evento; en las aves canoras las neuronas se disparan estrictamente por turno & Registro de neuronas identificadas del centro vocal superior del diamante mandarín durante el sueño y el canto: cada neurona que envía su salida al núcleo motor emite una sola ráfaga breve en un instante preciso del canto, y las neuronas se disparan en secuencia & \cite{Hahnloser2002} & medido \\
\bottomrule
\end{longtable}}
